\section{Conventions, flux and the oriented lens example}\label{app:dict}
For future-directed Lorentzian normal $n$, take the spatial orientation for which $c_1c_2c_3=-1$. Our symbol $P$ is the oriented integral of $e_i\wedge T^i$ and is the negative of the Lorentzian canonical momentum of $\theta$. In a spatial contortion convention $K_{ij}=w\epsilon_{ijk} e^k$ the axial parameter is $w=-s c/q^3$. The three-dimensional coefficient $+3w/2$ becomes $-sX/q$, so multiplying the proper-time Hamiltonian by $q$ gives the conformal shift $-sX$. The trace torsion from real $h$ is absent from the Hermitian spinor action. An unsymmetrized prescription or a reversal of normal and chirality changes the dictionary and the boundary projectors.

On a shell of unit-round absolute energy $k+3/2$, the positive spinors are $V_k\otimes V_{k+1}$ and the negative spinors $V_{k+1}\otimes V_k$ with $V_j=\mathrm{Sym}^j\mathbb C^2$ \cite{BarSpin,P08}. For our $L(7,2)$ choose the lift in the two $SU(2)$ factors with modular weights $(3,2)$: a weight pair $(u,v)$ survives when $3u+2v=0\bmod 7$. This explicitly distinguishes the representative from its mirror, which has $\eta_0=-4/7$. The same convention at $p=5$ uses $(2,1)$. Reversing the charge while also reversing orientation interchanges shell signs. An unaccompanied relabeling of filled and empty states changes the definition of number and cannot be made while retaining the original number operator.

There are two sorts of orientation at the boundary. Both spinor Fock spaces are transported with the same tangent frame to the common throat. Green's formula, by contrast, uses opposite outward normals at the two endpoints. With $\dd\tau=q\,\dd t$ the unitary transformation to half density leaves norm $\int|f|^2\dd\tau=\int q|f|^2\dd t$ in a normalized angular mode. The determinant calculation uses proper distance in the separated second-order operator; its transport exponent $\int b_s\,\dd\tau$ can be evaluated in conformal time. Interchanging these measures would change a heat trace.
